\scanpage{141}

(34)\qquad $\lambda\deleted{\unsure{(s)}}s$ est l'unique élément de $S$, dont l'étoile pour $A_\lambda=A_0^{\nu}$ est égale à $\mathrm{Et}_0(s)$.

$[$\uline{NB} Dans le cas de l'octaèdre, rappelons que $Z=\{1,\unsure{-|\underline{\circ}|\varepsilon}\}$, et que pour $s\in S$, $\unsure{\underline{\circ}|\underline{-}|\varepsilon}(s)$ peut se caractériser comme l'unique $s'\ne s$ \deleted{qui} \unsure{dont|qui} a même étoile que $s$. $]$ Ainsi on trouve un \uline{homomorphisme} \uline{injectif}
\[
Z\hookrightarrow\Lambda'_n\tag*{(35)}
\]
\margin{(cas différent du cas octaédral \deleted{\unsure{y|qui}} où on a $\Lambda'_n=\{1\}$, $Z=\{1,\underline{\circ}\}$)}
dont l'image est formée d'éléments $\nu$ satisfaisant la condition suivante

(36)\qquad L'ensemble \deleted{$\unsure{A'_0|A_0^{(\nu)}}=\{a\in\unsure{\mathfrak P|\mathcal P|P}_2(S)\mid\exists E\in\unsure{\mathcal E|\mathscr E|E}_0\}$} \deleted{$\unsure{\coprod_{E\in\mathcal E_0} A_E\text{ t.q. }a\in{}}$}
\[
A_0^{(\nu)}=\coprod_{E\in\mathcal E_0}A_E^{(\nu)}
\]
est l'ensemble des arêtes pour une structure de 2-polyèdre sur $S$ $[$ayant mêmes ``étoiles que $P_0$, et en fait, plus précisément, \uline{équivalente} à celle de $P_0$ au sens de (18)$]$.

Mais le simple fait déjà, \deleted{que} pour une structure de 2-polyèdre triangulé $P_0=(S,A_0,F_0)$ satisfaisant $[$To 1 et$]$ To 2,\inserted{ (To 5),} que l'ens $A_0^{(\nu)}\subset\mathfrak P_2(S)$ définisse sur $S$ une structure de graphe provenant, elle aussi, d'une telle structure triangulée -- et ceci \uline{pour tout $\nu\in\deleted{\unsure{Z'_n}}\Lambda'_n$} -- me paraît déjà une propriété remarquable et très forte, qui mérite d'être examinée

% slanted
\margin{cf To 7 plus bas !}

% archive sheet number: 140
